% Fichier engendré par tools/generate-tex.py à partir de 05-exponentielle-logarithme-trigonometrie.
% Les démonstrations en français sont transcrites du fichier Lean (skill transcribe-lean-proof).
\documentclass[11pt,a4paper]{article}

% Compile aussi bien avec pdflatex qu'avec xelatex ou tectonic.
\usepackage{iftex}
\ifPDFTeX
  \usepackage[T1]{fontenc}
  \usepackage[utf8]{inputenc}
\else
  \usepackage{fontspec}
\fi
\usepackage[french]{babel}
\usepackage{amsmath,amssymb,amsthm}
\usepackage[margin=2.5cm]{geometry}
% Les figures : TikZ, pour que leurs étiquettes soient composées dans les
% polices du document. Voir la skill `illustrate-theorem`.
\usepackage{tikz}
\usetikzlibrary{calc}

\usepackage[hidelinks]{hyperref}

\theoremstyle{definition}
\newtheorem{definition}{Définition}
\theoremstyle{plain}
\newtheorem{theoreme}{Théorème}
\newtheorem{lemme}[theoreme]{Lemme}

% Renvoi à la déclaration Lean, une ligne après la démonstration, aligné à droite.
\newcommand{\source}[2]{\par\smallskip\noindent\hfill{\footnotesize\href{#1}{\texttt{#2}}}\par\smallskip}

\title{Exponentielle, logarithme, trigonométrie}
\date{}

\begin{document}
\maketitle
% Les identifiants Lean en machine à écrire ne se coupent pas : on tolère des
% espacements plus lâches plutôt que des lignes qui débordent.
\sloppy

\section{ExponentielleLogarithmeTrigonometrie}

\noindent
Lycée \textemdash{} section \og{} Exponentielle, logarithme, trigonométrie \fg{}, niveau
première et terminale S.

\medskip\noindent
Une différence de construction traverse tout le chapitre. Le programme \emph{définit}
l'exponentielle comme l'unique fonction dérivable vérifiant $f' = f$ et $f(0) = 1$, et
admet que cette fonction existe. Mathlib procède autrement : $\exp$, $\sin$ et $\cos$ y
sont construits par leurs séries entières, et $\pi$ comme le double du premier zéro du
cosinus. Le premier énoncé rétablit le pont : l'exponentielle de la bibliothèque est
solution de l'équation différentielle, et le second démontre qu'elle en est la seule
solution \textemdash{} c'est la démonstration substantielle du chapitre, celle que le
programme admet.

\medskip\noindent
Les autres énoncés sont, pour l'essentiel, empruntés à Mathlib. Les transcriptions disent
ce que le résultat cité affirme et par quel argument il s'obtient, sans faire croire que le
travail est fait ici.

\subsection{Exponentielle : équation différentielle}

\begin{theoreme}
L'exponentielle est solution de l'équation différentielle du programme : elle est
dérivable, égale à sa dérivée, et vaut $1$ en zéro,
\[ \exp' = \exp, \qquad \exp(0) = 1 . \]
\end{theoreme}

\begin{proof}
Tout vient de la définition : $\exp$ est posée par la série
\[ \exp(x) = \sum_{n \ge 0} \frac{x^{n}}{n!} , \]
dont on dérive terme à terme pour obtenir $\exp' = \exp$, et dont le terme constant donne
$\exp(0) = 1$. Le fichier ne fait que les citer côte à côte : leur intérêt est de dire que
la fonction de Mathlib est bien celle que le programme définit.
\end{proof}
\source{https://github.com/Commutator-IO/learn-lean/blob/main/courses/02-lycee/05-exponentielle-logarithme-trigonometrie/ExponentielleLogarithmeTrigonometrie.lean\#L19}{ExponentielleLogarithmeTrigonometrie.lean\#L19}

\begin{theoreme}
Et c'est la seule : si $f$ est dérivable sur $\mathbb{R}$, vérifie $f'(x) = f(x)$ pour tout
$x$ et $f(0) = 1$, alors $f = \exp$.
\end{theoreme}

\begin{proof}
Posons $g(t) = f(t)\,e^{-t}$. La fonction $t \mapsto e^{-t}$ est dérivable de dérivée
$-e^{-t}$ \textemdash{} c'est la dérivée d'une composée, l'exponentielle étant composée avec
$t \mapsto -t$ dont la dérivée vaut $-1$. La règle du produit donne alors, en tout point
$y$,
\[ g'(y) = f'(y)\,e^{-y} + f(y)\,\bigl(-e^{-y}\bigr)
        = f(y)\,e^{-y} - f(y)\,e^{-y} = 0 , \]
où l'on a remplacé $f'(y)$ par $f(y)$ selon l'hypothèse.

Une fonction dérivable sur $\mathbb{R}$ de dérivée identiquement nulle est constante
(conséquence du théorème des accroissements finis, démontrée dans la bibliothèque). Donc
$g(x) = g(0)$ pour tout $x$, et
\[ f(x)\,e^{-x} = f(0)\,e^{0} = 1 . \]
Comme $e^{-x} = 1/e^{x}$ et que l'exponentielle ne s'annule pas, on peut multiplier par
$e^{x}$ : $f(x) = e^{x}$.
\end{proof}
\source{https://github.com/Commutator-IO/learn-lean/blob/main/courses/02-lycee/05-exponentielle-logarithme-trigonometrie/ExponentielleLogarithmeTrigonometrie.lean\#L25}{ExponentielleLogarithmeTrigonometrie.lean\#L25}

\subsection{Propriétés algébriques de l'exponentielle}

\begin{theoreme}
L'exponentielle transforme les sommes en produits :
\[ \exp(a + b) = \exp(a)\exp(b), \qquad \exp(-a) = \frac{1}{\exp(a)}, \]
\[ \exp(na) = \exp(a)^{n} \quad (n \in \mathbb{N}) . \]
\end{theoreme}

\begin{proof}
La première égalité vient du produit de Cauchy des deux séries. En écrivant $\exp(a) = \sum_n a^n/n!$ et de même pour $b$, le produit des deux séries se
réordonne en
\[ \sum_{n} \frac{1}{n!} \sum_{k=0}^{n} \binom{n}{k} a^{k} b^{n-k}
   = \sum_{n} \frac{(a+b)^n}{n!} = \exp(a + b) , \]
la seconde égalité étant la formule du binôme. Le réordonnancement est licite parce que les
deux séries convergent absolument.

Les deux autres s'en déduisent. En prenant $b = -a$, il vient $\exp(a)\exp(-a) = \exp(0) = 1$,
donc $\exp(-a) = 1/\exp(a)$ ; et une récurrence sur $n$, partant de $\exp(0) = 1$, donne
$\exp(na) = \exp(a)^{n}$.

\medskip\noindent
C'est cette propriété qui vaut à $\exp$ son nom : elle échange la structure additive de
$\mathbb{R}$ et la structure multiplicative de $\mathbb{R}_{+}^{*}$.
\end{proof}
\source{https://github.com/Commutator-IO/learn-lean/blob/main/courses/02-lycee/05-exponentielle-logarithme-trigonometrie/ExponentielleLogarithmeTrigonometrie.lean\#L45}{ExponentielleLogarithmeTrigonometrie.lean\#L45}

\begin{theoreme}
L'exponentielle est strictement positive et strictement croissante :
\[ \exp(x) > 0 \quad \text{pour tout } x, \qquad
   x < y \implies \exp(x) < \exp(y) . \]
\end{theoreme}

\begin{proof}
La positivité tient à l'identité
$\exp(x) = \exp(x/2)^{2}$, qui donne un carré, non nul puisque l'exponentielle ne s'annule
jamais. La croissance stricte s'en déduit : pour $x < y$,
\[ \exp(y) - \exp(x) = \exp(x)\bigl(\exp(y - x) - 1\bigr) > 0 , \]
le second facteur étant strictement positif dès que $y - x > 0$.
\end{proof}
\source{https://github.com/Commutator-IO/learn-lean/blob/main/courses/02-lycee/05-exponentielle-logarithme-trigonometrie/ExponentielleLogarithmeTrigonometrie.lean\#L54}{ExponentielleLogarithmeTrigonometrie.lean\#L54}

\subsection{Limites de l'exponentielle}

\begin{theoreme}
\[ e^{x} \xrightarrow[x \to +\infty]{} +\infty , \qquad
   e^{x} \xrightarrow[x \to -\infty]{} 0 . \]
\end{theoreme}

\begin{proof}
La première limite vient de la minoration $\exp(x) \ge 1 + x$, valable pour tout $x$ ; la seconde s'en déduit par
$\exp(-x) = 1/\exp(x)$, l'inverse d'une quantité tendant vers $+\infty$ tendant vers zéro.
\end{proof}
\source{https://github.com/Commutator-IO/learn-lean/blob/main/courses/02-lycee/05-exponentielle-logarithme-trigonometrie/ExponentielleLogarithmeTrigonometrie.lean\#L61}{ExponentielleLogarithmeTrigonometrie.lean\#L61}

\subsection{Logarithme népérien}

\begin{theoreme}
Le logarithme est la fonction réciproque de l'exponentielle :
\[ \ln(e^{x}) = x \ \text{ pour tout } x, \qquad
   e^{\ln x} = x \ \text{ pour } x > 0 . \]
\end{theoreme}

\begin{proof}
Les deux identités disent que $\ln$ est la réciproque de $\exp$ sur $\mathbb{R}_{+}^{*}$ :
c'est ainsi qu'elle est définie, et il n'y a rien d'autre à établir.

\medskip\noindent
L'asymétrie des hypothèses n'est pas un oubli : $\ln(e^{x}) = x$ vaut pour tout réel
puisque $e^{x} > 0$, tandis que $e^{\ln x} = x$ réclame $x > 0$. Le logarithme de Lean est
pourtant une fonction totale \textemdash{} la bibliothèque pose $\ln x = \ln|x|$ pour
$x < 0$ et $\ln 0 = 0$ \textemdash{} mais ces valeurs de remplissage ne sont pas des
mathématiques : elles rendent la fonction totale, rien de plus, et l'hypothèse $x > 0$ reste
celle de l'énoncé.
\end{proof}
\source{https://github.com/Commutator-IO/learn-lean/blob/main/courses/02-lycee/05-exponentielle-logarithme-trigonometrie/ExponentielleLogarithmeTrigonometrie.lean\#L69}{ExponentielleLogarithmeTrigonometrie.lean\#L69}

\begin{figure}[h]
\centering
% Construction : les courbes de exp et de ln, symétriques l'une de l'autre par
% rapport à la première bissectrice, tracée en pointillés.
\begin{tikzpicture}[scale=0.8]
  \draw[->] (-2.2, 0) -- (3.6, 0) node[right] {$x$};
  \draw[->] (0, -2.2) -- (0, 3.6) node[above] {$y$};
  \draw[dashed] (-1.8, -1.8) -- (3.2, 3.2);
  \draw[thick, domain=-2:1.15, smooth, variable=\x] plot ({\x}, {exp(\x)});
  \draw[thick, domain=0.14:3.15, smooth, variable=\x] plot ({\x}, {ln(\x)});
  \node[above left] at (1.15, 3.16) {$\exp$};
  \node[below right] at (3.15, 1.15) {$\ln$};
  \node[above left] at (3.2, 3.2) {$y = x$};
\end{tikzpicture}

\caption{Les courbes de $\exp$ et de $\ln$ sont symétriques par rapport à la première bissectrice.}
\end{figure}

\begin{theoreme}
Pour $a > 0$ et $b > 0$, le logarithme transforme les produits en sommes :
\[ \ln(ab) = \ln a + \ln b, \qquad \ln\frac{a}{b} = \ln a - \ln b, \]
\[ \ln(a^{n}) = n \ln a \quad (n \in \mathbb{N}), \qquad
   \ln\sqrt{a} = \tfrac{1}{2}\ln a . \]
\end{theoreme}

\begin{proof}
Toutes découlent de la première par la réciprocité : $\ln(ab) = \ln a + \ln b$ est la
traduction, par $\ln$, de $e^{u+v} = e^{u}e^{v}$.

\medskip\noindent
Sur les hypothèses, Mathlib demande moins : $a \ne 0$ et $b \ne 0$ suffisent, parce que le
logarithme de Lean d'un négatif vaut celui de sa valeur absolue. La transcription retient
celles du lycée, $a > 0$ et $b > 0$, qui sont celles sous lesquelles l'énoncé a un sens.
\end{proof}
\source{https://github.com/Commutator-IO/learn-lean/blob/main/courses/02-lycee/05-exponentielle-logarithme-trigonometrie/ExponentielleLogarithmeTrigonometrie.lean\#L75}{ExponentielleLogarithmeTrigonometrie.lean\#L75}

\begin{theoreme}
Le logarithme est dérivable en tout point non nul, de dérivée
\[ (\ln x)' = \frac{1}{x} , \]
et il est strictement croissant sur $]0\,;+\infty[$.
\end{theoreme}

\begin{proof}
La dérivée s'obtient par dérivation de la réciproque : en dérivant $e^{\ln x} = x$ on
trouve $e^{\ln x} \times (\ln x)' = 1$, soit $x\,(\ln x)' = 1$. Mathlib l'écrit $x^{-1}$,
que l'on récrit $1/x$.

La croissance stricte tient à ce que $\ln$ est la réciproque d'une fonction strictement
croissante : elle conserve donc les inégalités strictes entre réels strictement positifs.
La transcription l'applique à deux points $x < y$ de $]0\,;+\infty[$ ; l'appartenance de $y$ à
l'intervalle ne sert pas, seule celle de $x$ est requise.
\end{proof}
\source{https://github.com/Commutator-IO/learn-lean/blob/main/courses/02-lycee/05-exponentielle-logarithme-trigonometrie/ExponentielleLogarithmeTrigonometrie.lean\#L85}{ExponentielleLogarithmeTrigonometrie.lean\#L85}

\begin{theoreme}
\[ \ln x \xrightarrow[\substack{x \to 0 \\ x > 0}]{} -\infty , \qquad
   \ln x \xrightarrow[x \to +\infty]{} +\infty . \]
\end{theoreme}

\begin{proof}
Elles se déduisent des limites de l'exponentielle par réciprocité.

Le logarithme est défini comme la réciproque de $\exp$, laquelle est une bijection
strictement croissante de $\mathbb{R}$ sur $\mathbb{R}_{+}^{*}$. Une bijection croissante et
sa réciproque échangent les extrémités : comme $\exp$ envoie $-\infty$ sur $0^{+}$ et
$+\infty$ sur $+\infty$, le logarithme envoie $0^{+}$ sur $-\infty$ et $+\infty$ sur
$+\infty$.

Concrètement, pour la première : soit $A > 0$ ; pour $0 < x < \exp(-A)$, la croissance du
logarithme donne $\ln x < -A$. Pour la seconde : pour $x > \exp(A)$, on a $\ln x > A$.
\end{proof}
\source{https://github.com/Commutator-IO/learn-lean/blob/main/courses/02-lycee/05-exponentielle-logarithme-trigonometrie/ExponentielleLogarithmeTrigonometrie.lean\#L93}{ExponentielleLogarithmeTrigonometrie.lean\#L93}

\subsection{Cercle trigonométrique}

\begin{theoreme}
Le point de coordonnées $(\cos x, \sin x)$ est sur le cercle de centre l'origine et de
rayon $1$ :
\[ \cos^{2} x + \sin^{2} x = 1 , \]
et les deux coordonnées sont donc comprises entre $-1$ et $1$.
\end{theoreme}

\begin{proof}
Les trois faits sont cités à la bibliothèque. L'identité de Pythagore s'y obtient à partir
des séries définissant $\cos$ et $\sin$ ; les deux majorations $|\cos x| \le 1$ et
$|\sin x| \le 1$ en découlent, un carré ne pouvant excéder la somme.

\medskip\noindent
Ce que le fichier ne formalise pas : que le paramètre $x$ mesure la \emph{longueur} de
l'arc parcouru sur ce cercle, c'est-à-dire la définition même du radian. Établir ce point
demande la longueur d'un arc, donc une intégrale, absente de ce chapitre. Le cercle
trigonométrique y est donc présent par son équation, non par sa géométrie.
\end{proof}
\source{https://github.com/Commutator-IO/learn-lean/blob/main/courses/02-lycee/05-exponentielle-logarithme-trigonometrie/ExponentielleLogarithmeTrigonometrie.lean\#L102}{ExponentielleLogarithmeTrigonometrie.lean\#L102}

\begin{theoreme}
Un tour complet vaut $2\pi$ radians : le cosinus et le sinus sont $2\pi$-périodiques,
\[ \cos(x + 2\pi) = \cos x, \qquad \sin(x + 2\pi) = \sin x . \]
\end{theoreme}

\begin{proof}
Dans Mathlib, $\pi$ est \emph{défini} comme le double du premier zéro du cosinus sur
$[0\,;2]$. La périodicité s'en déduit par les formules d'addition, en deux temps.

De $\cos(\pi/2) = 0$ et $\cos^2 + \sin^2 = 1$ on tire $\sin(\pi/2) = 1$, le sinus étant
positif sur $[0\,;\pi/2]$. Les formules d'addition donnent alors
\[ \cos \pi = \cos^2(\pi/2) - \sin^2(\pi/2) = -1 , \qquad
   \sin \pi = 2\sin(\pi/2)\cos(\pi/2) = 0 , \]
puis, de la même façon, $\cos 2\pi = 1$ et $\sin 2\pi = 0$. Il ne reste qu'à développer :
\[ \cos(x + 2\pi) = \cos x \cos 2\pi - \sin x \sin 2\pi = \cos x , \]
et de même pour le sinus.

La définition scolaire \textemdash{} $\pi$ est le rapport de la circonférence au diamètre
\textemdash{} n'est pas celle-là, et le fichier ne démontre pas qu'elles coïncident.
\end{proof}
\source{https://github.com/Commutator-IO/learn-lean/blob/main/courses/02-lycee/05-exponentielle-logarithme-trigonometrie/ExponentielleLogarithmeTrigonometrie.lean\#L107}{ExponentielleLogarithmeTrigonometrie.lean\#L107}

\subsection{Valeurs remarquables et angles associés}

\begin{theoreme}
Valeurs remarquables sur le premier quadrant :
\[ \cos 0 = 1, \quad \cos\frac{\pi}{6} = \frac{\sqrt{3}}{2}, \quad
   \cos\frac{\pi}{4} = \frac{\sqrt{2}}{2}, \quad \cos\frac{\pi}{3} = \frac{1}{2}, \quad
   \cos\frac{\pi}{2} = 0, \]
\[ \sin 0 = 0, \quad \sin\frac{\pi}{6} = \frac{1}{2}, \quad
   \sin\frac{\pi}{4} = \frac{\sqrt{2}}{2}, \quad \sin\frac{\pi}{3} = \frac{\sqrt{3}}{2},
   \quad \sin\frac{\pi}{2} = 1 . \]
\end{theoreme}

\begin{proof}
Ces dix valeurs ne se calculent pas à partir des séries, mais des relations remarquables : $\cos(\pi/2) = 0$ tient à la
définition de $\pi$, $\cos(\pi/3) = 1/2$ vient de la formule de duplication appliquée à
$\pi/3$, et $\cos(\pi/4) = \sqrt{2}/2$ de $\cos^{2} + \sin^{2} = 1$ jointe à
$\cos(\pi/4) = \sin(\pi/4)$.
\end{proof}
\source{https://github.com/Commutator-IO/learn-lean/blob/main/courses/02-lycee/05-exponentielle-logarithme-trigonometrie/ExponentielleLogarithmeTrigonometrie.lean\#L114}{ExponentielleLogarithmeTrigonometrie.lean\#L114}

\begin{theoreme}
Angles associés, qui se lisent par symétrie sur le cercle :
\[ \cos(-x) = \cos x, \qquad \sin(-x) = -\sin x, \]
\[ \cos(\pi - x) = -\cos x, \qquad \sin(\pi - x) = \sin x, \]
\[ \cos\Bigl(\frac{\pi}{2} - x\Bigr) = \sin x, \qquad
   \sin\Bigl(\frac{\pi}{2} - x\Bigr) = \cos x . \]
\end{theoreme}

\begin{proof}
Toutes se déduisent des formules d'addition jointes aux valeurs remarquables.

La parité du cosinus et l'imparité du sinus se lisent sur les séries, où le cosinus n'a que
des puissances paires et le sinus que des impaires. Les quatre autres s'obtiennent en
développant :
\[ \cos(\pi - x) = \cos\pi\cos x + \sin\pi\sin x = -\cos x , \]
\[ \sin(\pi - x) = \sin\pi\cos x - \cos\pi\sin x = \sin x , \]
avec $\cos\pi = -1$ et $\sin\pi = 0$ ; et de même en $\pi/2$, où $\cos(\pi/2) = 0$ et
$\sin(\pi/2) = 1$ donnent
\[ \cos(\pi/2 - x) = \sin x , \qquad \sin(\pi/2 - x) = \cos x . \]

Géométriquement, ces trois couples sont trois symétries du cercle trigonométrique :
l'axe des abscisses, l'axe des ordonnées, et la première bissectrice, qui échange les deux
coordonnées.
\end{proof}
\source{https://github.com/Commutator-IO/learn-lean/blob/main/courses/02-lycee/05-exponentielle-logarithme-trigonometrie/ExponentielleLogarithmeTrigonometrie.lean\#L125}{ExponentielleLogarithmeTrigonometrie.lean\#L125}

\subsection{Formules d'addition et de duplication}

\begin{theoreme}
\[ \cos(a + b) = \cos a \cos b - \sin a \sin b, \qquad
   \sin(a + b) = \sin a \cos b + \cos a \sin b, \]
\[ \cos(2a) = \cos^{2} a - \sin^{2} a, \qquad \sin(2a) = 2 \sin a \cos a . \]
\end{theoreme}

\begin{proof}
Le chemin le plus court passe par l'exponentielle complexe et l'identité $e^{i\theta} = \cos\theta +
i\sin\theta$ : la relation $e^{ia}e^{ib} = e^{i(a+b)}$ s'écrit
\[ (\cos a + i\sin a)(\cos b + i\sin b)
   = \cos(a+b) + i\sin(a+b) , \]
et le développement du membre de gauche donne
\[ (\cos a\cos b - \sin a\sin b) + i\,(\sin a\cos b + \cos a\sin b) . \]
L'identification des parties réelle et imaginaire rend d'un coup les deux formules
d'addition.

Les formules de duplication en sont le cas $b = a$.

\medskip\noindent
Le raisonnement n'est pas circulaire : l'exponentielle complexe est définie par sa série,
indépendamment de $\cos$ et $\sin$, et l'identité d'Euler se démontre en séparant les termes
pairs et impairs de cette série.
\end{proof}
\source{https://github.com/Commutator-IO/learn-lean/blob/main/courses/02-lycee/05-exponentielle-logarithme-trigonometrie/ExponentielleLogarithmeTrigonometrie.lean\#L135}{ExponentielleLogarithmeTrigonometrie.lean\#L135}

\subsection{Résolution des équations trigonométriques}

\begin{theoreme}
Pour $-1 \le a \le 1$, l'équation $\cos x = a$ admet $\arccos a$ pour solution, et ses
solutions sont exactement
\[ x = 2k\pi + \arccos a \quad \text{ou} \quad x = 2k\pi - \arccos a,
   \qquad k \in \mathbb{Z} . \]
\end{theoreme}

\begin{proof}
La bibliothèque fournit deux ingrédients : $\cos(\arccos a) = a$ dès que
$-1 \le a \le 1$, et la caractérisation
\[ \cos u = \cos v \iff \exists k \in \mathbb{Z},\
   v = 2k\pi + u \ \text{ ou } \ v = 2k\pi - u . \]

Dans un sens : si $\cos x = a$, alors $\cos(\arccos a) = a = \cos x$, et la
caractérisation appliquée à $u = \arccos a$, $v = x$ donne la forme annoncée.

Dans l'autre : si $x$ a l'une des deux formes, la caractérisation donne
$\cos(\arccos a) = \cos x$, et le membre de gauche vaut $a$.

\medskip\noindent
Les deux familles de solutions sont les deux points du cercle de même abscisse $a$,
symétriques par rapport à l'axe des abscisses.
\end{proof}
\source{https://github.com/Commutator-IO/learn-lean/blob/main/courses/02-lycee/05-exponentielle-logarithme-trigonometrie/ExponentielleLogarithmeTrigonometrie.lean\#L146}{ExponentielleLogarithmeTrigonometrie.lean\#L146}

\begin{theoreme}
Pour $-1 \le a \le 1$, l'équation $\sin x = a$ admet $\arcsin a$ pour solution, et ses
solutions sont exactement
\[ x = 2k\pi + \arcsin a \quad \text{ou} \quad x = (2k+1)\pi - \arcsin a,
   \qquad k \in \mathbb{Z} . \]
\end{theoreme}

\begin{proof}
Même démonstration, avec $\sin(\arcsin a) = a$ et la caractérisation
\[ \sin u = \sin v \iff \exists k \in \mathbb{Z},\
   v = 2k\pi + u \ \text{ ou } \ v = (2k+1)\pi - u . \]
La seconde famille correspond au point symétrique par rapport à l'axe des ordonnées : même
ordonnée $a$, abscisse opposée.
\end{proof}
\source{https://github.com/Commutator-IO/learn-lean/blob/main/courses/02-lycee/05-exponentielle-logarithme-trigonometrie/ExponentielleLogarithmeTrigonometrie.lean\#L162}{ExponentielleLogarithmeTrigonometrie.lean\#L162}

\begin{theoreme}
Une équation trigonométrique a donc une infinité de solutions : à toute solution on peut
ajouter un nombre entier de tours,
\[ \cos(x + 2k\pi) = \cos x, \qquad \sin(x + 2k\pi) = \sin x, \qquad k \in \mathbb{Z} . \]
\end{theoreme}

\begin{proof}
La bibliothèque énonce ces deux identités sous la forme $\cos(x + k \times 2\pi) = \cos x$.
La transcription se ramène à cette écriture par
\[ 2k\pi = k \times (2\pi) , \]
simple réassociation des facteurs.
\end{proof}
\source{https://github.com/Commutator-IO/learn-lean/blob/main/courses/02-lycee/05-exponentielle-logarithme-trigonometrie/ExponentielleLogarithmeTrigonometrie.lean\#L177}{ExponentielleLogarithmeTrigonometrie.lean\#L177}

\subsection{Dérivées du sinus et du cosinus}

\begin{theoreme}
\[ \sin' = \cos, \qquad \cos' = -\sin . \]
\end{theoreme}

\begin{proof}
Les deux dérivées s'obtiennent en dérivant terme à terme les séries entières qui définissent
$\sin$ et $\cos$ dans la bibliothèque — ce qui est licite à l'intérieur du disque de
convergence, ici le plan tout entier.

En dérivant
\[ \sin x = x - \frac{x^3}{3!} + \frac{x^5}{5!} - \cdots \]
on obtient
\[ 1 - \frac{x^2}{2!} + \frac{x^4}{4!} - \cdots = \cos x , \]
et la même opération sur la série du cosinus rend $-\sin$, le premier terme constant
disparaissant et les signes se décalant.

C'est un ordre de présentation inverse de celui du lycée, qui pose $\sin$ et $\cos$
géométriquement et admet leurs dérivées.
\end{proof}
\source{https://github.com/Commutator-IO/learn-lean/blob/main/courses/02-lycee/05-exponentielle-logarithme-trigonometrie/ExponentielleLogarithmeTrigonometrie.lean\#L187}{ExponentielleLogarithmeTrigonometrie.lean\#L187}

\begin{theoreme}
\[ \frac{\sin x}{x} \xrightarrow[\substack{x \to 0 \\ x \ne 0}]{} 1 . \]
\end{theoreme}

\begin{proof}
Cette limite est le nombre dérivé du sinus en zéro. On part de $\sin'(0) = \cos 0 = 1$ et
on l'écrit comme limite du taux d'accroissement : la pente du sinus entre $0$ et $x$ vaut
\[ \frac{\sin x - \sin 0}{x - 0} = \frac{\sin x}{x} , \]
puisque $\sin 0 = 0$. La dérivabilité en zéro dit exactement que ce quotient tend vers $1$
lorsque $x$ tend vers zéro en restant distinct de zéro.

\medskip\noindent
Au lycée, l'ordre est inverse : on établit d'abord $\sin x / x \to 1$, par un encadrement
géométrique, pour en déduire $\sin' = \cos$. Ici il n'y a pas de cercle vicieux, parce que
la bibliothèque démontre $\sin' = \cos$ à partir des séries, sans passer par cette limite.
\end{proof}
\source{https://github.com/Commutator-IO/learn-lean/blob/main/courses/02-lycee/05-exponentielle-logarithme-trigonometrie/ExponentielleLogarithmeTrigonometrie.lean\#L192}{ExponentielleLogarithmeTrigonometrie.lean\#L192}

\vfill
\noindent\rule{0.35\linewidth}{0.4pt}\par
\noindent{\footnotesize Leonardo de Moura et Sebastian Ullrich,
\emph{The Lean~4 Theorem Prover and Programming Language},
\emph{Automated Deduction — CADE~28}, Springer, 2021, p.~625--635,
\href{https://doi.org/10.1007/978-3-030-79876-5_37}{doi:10.1007/978-3-030-79876-5\_37}.\par}

\end{document}
